\section{Methodology}
\label{sec:methodology}

\subsection{Overview}

The approach has four components: rolling IC estimation with exponential decay, regime detection and conditional multipliers, walk-forward validation, and risk control.

\subsection{Factor Library}

We designed a 25-factor library across five categories. This study implements 10 core factors. The full 25-factor framework is in Appendix~\ref{sec:factor_appendix}.

\begin{table}[h]
\centering
\caption{Core Factors Implemented (10 of 25)}
\label{tab:factors}
\begin{tabular}{llc}
\toprule
Category & Factors & Count \\
\midrule
Momentum & return\_20d, return\_60d, return\_120d, ma\_return\_20d, momentum\_12m & 5 \\
Reversal & reversal\_5d, reversal\_20d & 2 \\
Volatility & volatility\_20d, volatility\_60d & 2 \\
Liquidity & volume\_ratio & 1 \\
\midrule
Total & & 10 \\
\bottomrule
\end{tabular}
\end{table}

We chose these 10 for data availability and computational simplicity. Momentum factors measure past cumulative returns. Reversal factors capture short-term mean reversion. Volatility factors measure risk characteristics. All factors are winsorized at 1st/99th percentiles and standardized cross-sectionally.

\subsection{Rolling IC Estimation}

The Information Coefficient for factor $f$ at time $t$ is the Spearman rank correlation between factor values and forward returns:

\begin{equation}
IC_{f,t} = \rho_{\text{Spearman}}(\mathbf{f}_t, \mathbf{r}_{t+1:t+N})
\end{equation}

We use $N=20$ trading days as the prediction horizon. To smooth noise while tracking recent effectiveness, we compute decay-weighted IC mean:

\begin{equation}
\overline{IC}_{f,t} = \frac{\sum_{i=0}^{W-1} \gamma^i \cdot IC_{f,t-i}}{\sum_{i=0}^{W-1} \gamma^i}
\end{equation}

\begin{equation}
ICIR_{f,t} = \frac{\overline{IC}_{f,t}}{\sigma_{IC,f,t}}
\end{equation}

where $W = 60$ trading days is the lookback window and $\gamma \in (0,1)$ is the decay factor.

\subsection{Weight Computation}

We exclude factors with negative average IC:

\begin{equation}
w_{f,t}^{\text{base}} = \begin{cases}
\min\left(\frac{\overline{IC}_{f,t}}{\sigma_{IC,f,t}} \cdot \eta, w_{\text{max}}\right) & \text{if } \overline{IC}_{f,t} > 0 \\
0 & \text{if } \overline{IC}_{f,t} \leq 0
\end{cases}
\end{equation}

where $\eta = 0.1$ scales the weights and $w_{\text{max}} = 0.15$ caps individual factors. Only factors with positive predictive power get positive weights.

Weights are normalized:

\begin{equation}
w_{f,t} = \frac{w_{f,t}^{\text{base}}}{\sum_{f'} w_{f',t}^{\text{base}}}
\end{equation}

\subsection{Regime Detection}

We classify market conditions into Bull, Bear, or Oscillation using trend and volatility indicators. Bull: price above 60-day MA, volatility below 20\% annualized, positive 20-day momentum. Bear: price below 60-day MA, volatility above 25\% annualized, negative momentum. Oscillation: neither condition met.

Regime-specific multipliers adjust factor weights based on documented effectiveness patterns~\cite{arnott2016timing,blitz2020factormomentum}:

\begin{table}[h]
\centering
\caption{Regime-Specific Factor Multipliers}
\label{tab:multipliers}
\begin{tabular}{lccc}
\toprule
Factor Category & Bull & Bear & Oscillation \\
\midrule
Momentum & 1.3$\times$ & 0.7$\times$ & 1.0$\times$ \\
Quality & 0.8$\times$ & 1.2$\times$ & 1.0$\times$ \\
Liquidity & 0.9$\times$ & 1.1$\times$ & 1.0$\times$ \\
Volatility & 1.0$\times$ & 1.1$\times$ & 1.0$\times$ \\
Sentiment & 1.0$\times$ & 1.0$\times$ & 1.0$\times$ \\
\bottomrule
\end{tabular}
\end{table}

The adjusted weight is:
\begin{equation}
w_{f,t}^{\text{adjusted}} = w_{f,t} \cdot m_{\text{regime},f}
\end{equation}

\subsection{Risk Control}

Maximum drawdown constraint: reduce position 50\% if drawdown exceeds 25\%, liquidate if exceeds 30\%. Volatility targeting: scale positions to maintain 15\% annualized volatility. Bear regime deleveraging: reduce positions 50\% automatically.

\subsection{Walk-Forward Validation}

Training window: 12 months (250 trading days). Test window: 3 months (60 trading days). Monthly rebalancing. This prevents look-ahead bias. Weights at any point depend only on information available then.

\subsection{Portfolio Construction}

Monthly rebalancing procedure: compute IC estimates (60-day window), detect regime, calculate adjusted weights, apply risk constraints, rank stocks by composite score, select top 30 stocks, equal-weight within portfolio.

Transaction costs: 0.03\% commission, 0.1\% stamp duty (sell side), 0.01\% slippage. These reflect actual Chinese A-share conditions.

\subsection{Parameter Sensitivity}

\begin{table}[h]
\centering
\caption{Parameter Sensitivity Analysis}
\label{tab:sensitivity}
\begin{tabular}{lcccc}
\toprule
Parameter & Values Tested & Sharpe Range & Max DD Range & Optimal \\
\midrule
Lookback Window $W$ & 30, 60, 90 days & [0.15, 0.18] & [28\%, 35\%] & 60 \\
Decay Factor $\gamma$ & 0.90, 0.95, 0.98 & [0.16, 0.18] & [29\%, 33\%] & 0.95 \\
Max Weight $w_{\max}$ & 0.10, 0.15, 0.20 & [0.17, 0.18] & [28\%, 32\%] & 0.15 \\
DD Limit $D_{\max}$ & 0.25, 0.30, 0.35 & [0.16, 0.19] & [25\%, 35\%] & 0.30 \\
\bottomrule
\end{tabular}
\end{table}

Performance was stable across parameter ranges. Optimal values: $W=60$, $\gamma=0.95$, $w_{\max}=0.15$, $D_{\max}=0.30$.

\subsection{Baseline Comparisons}

We tested against: equal-weighted factors (all 10 factors same weight), static IC weights (full-period IC, no updating), IC weights without regime adjustment.